\documentclass[11pt]{article}
\usepackage[margin=1in]{geometry}
\usepackage[T1]{fontenc}
\usepackage{lmodern}
\usepackage{microtype}
\usepackage{amsmath,amssymb}
\usepackage{booktabs,longtable,array}
\usepackage{enumitem}
\usepackage{xcolor}
\usepackage{hyperref}
\usepackage{fancyhdr}
\usepackage{listings}

\definecolor{ink}{HTML}{101814}
\definecolor{acid}{HTML}{79A11E}
\definecolor{paper}{HTML}{F3F0E7}
\definecolor{blue}{HTML}{3D7EA6}
\hypersetup{
  colorlinks=true,
  linkcolor=ink,
  urlcolor=blue,
  pdftitle={Profit-Yield Regime Monitor},
  pdfauthor={Ishva Mehta}
}
\pagestyle{fancy}
\fancyhf{}
\lhead{\textsc{Profit-Yield Regime Monitor}}
\rhead{\textsc{Technical Note}}
\cfoot{\thepage}
\setlength{\headheight}{14pt}
\setlength{\parindent}{0pt}
\setlength{\parskip}{0.65em}
\setlist{leftmargin=1.5em,itemsep=0.35em}
\renewcommand{\arraystretch}{1.2}

\lstdefinestyle{code}{
  language=Python,
  backgroundcolor=\color{paper},
  basicstyle=\ttfamily\small,
  keywordstyle=\color{blue}\bfseries,
  commentstyle=\color{acid},
  numbers=left,
  numberstyle=\tiny\color{gray},
  frame=single,
  breaklines=true,
  showstringspaces=false
}

\title{\vspace{-2em}
{\Large\textsc{U.S. Capital Markets}}\\[0.7em]
{\huge\bfseries Profit-Yield Regime Monitor}\\[0.5em]
{\large Methodology, Implementation, and Interpretation}}
\author{Ishva Mehta}
\date{\today}

\begin{document}
\maketitle

\begin{abstract}
The Profit-Yield Regime Monitor studies the interaction between U.S. corporate
profit growth and long-term Treasury yields. It uses public data from the U.S.
Bureau of Economic Analysis and the Federal Reserve Bank of St. Louis to
classify the market environment into one of four regimes. The project is
built as a small, reproducible research exercise. The data sources, signal
definitions, classification rule, and limitations are stated explicitly.
\end{abstract}

\tableofcontents
\newpage

\section{Research Question}

The monitor addresses a simple question:

\begin{quote}
\emph{What does the joint direction of corporate profit growth and long-term
government bond yields imply about the current market regime?}
\end{quote}

Rising profits are usually associated with stronger economic activity and
support for equities. Rising long-term yields can signal stronger nominal
growth, higher inflation risk, tighter financial conditions, or a higher term
premium. The two signals should therefore be interpreted together.

The project does not forecast security returns. It organizes two observable
signals into a weekly market note.

\section{Why Use a Public Earnings Proxy}

\subsection{Target concept}

Analyst earnings revisions are useful because they measure changes in
forward-looking expectations. High-quality historical estimate data are
usually supplied by commercial vendors. A public GitHub project should not
assume access to proprietary data.

\subsection{Selected proxy}

The monitor uses Corporate Profits After Tax from the Bureau of Economic
Analysis, available through FRED as series \texttt{CP}. This series is:

\begin{itemize}
  \item public;
  \item quarterly;
  \item seasonally adjusted at an annual rate;
  \item available from 1947; and
  \item revised as national accounts are updated.
\end{itemize}

The proxy measures realized, economy-wide profits. It is not equivalent to
forward-looking S\&P 500 analyst estimates. The distinction is shown directly
in the application and README.

\subsection{Trade-off}

\begin{longtable}{p{0.29\textwidth}p{0.29\textwidth}p{0.29\textwidth}}
\toprule
& \textbf{Analyst estimates} & \textbf{BEA proxy} \\
\midrule
\endhead
Orientation & Forward-looking & Realized \\
Coverage & Listed index companies & U.S. corporations \\
Frequency & Frequent revisions & Quarterly \\
History & Vendor-dependent & Public history from 1947 \\
Revisions & Forecast changes & National-account revisions \\
Use in this project & Target concept & Implemented signal \\
\bottomrule
\end{longtable}

\section{Data}

\subsection{Series}

\begin{table}[h]
\centering
\begin{tabular}{lll}
\toprule
\textbf{Series} & \textbf{Description} & \textbf{Frequency} \\
\midrule
\texttt{CP} & Corporate profits after tax & Quarterly \\
\texttt{DGS10} & 10-year Treasury yield & Daily \\
\texttt{DGS2} & 2-year Treasury yield & Daily \\
\texttt{T10YIE} & 10-year breakeven inflation & Daily \\
\bottomrule
\end{tabular}
\caption{Public FRED series used by the monitor}
\end{table}

The app downloads each series from FRED's public CSV endpoint. It does
not require an API key.

\subsection{Frequency alignment}

Treasury and breakeven data are converted to weekly Friday observations using
the last available value. Quarterly profit observations are carried forward
until a new observation becomes available. This creates a weekly table while
preserving the quarterly frequency of the profit signal.

The resulting weekly date does not mean profit data were updated that week. It
means the latest available quarterly profit observation was used for that
week's classification.

\section{Signal Construction}

\subsection{Profit growth}

Year-over-year corporate profit growth is:

\[
g_t = 100\left(\frac{CP_t}{CP_{t-4}}-1\right).
\]

Profit acceleration is the change in year-over-year growth:

\[
a_t = g_t-g_{t-1}.
\]

The profit signal is:

\[
P_t =
\begin{cases}
\text{accelerating}, & a_t>0,\\
\text{decelerating}, & a_t\leq0.
\end{cases}
\]

This signal measures whether the growth rate of profits is improving, not
whether the level of profits is positive.

\subsection{Yield trend}

Let $y_t^{10}$ be the 10-year Treasury yield and let $L$ be the selected
lookback in weeks. The change is:

\[
\Delta y_t^{10}=100\left(y_t^{10}-y_{t-L}^{10}\right)
\]

in basis points. The yield signal is:

\[
Y_t =
\begin{cases}
\text{rising}, & \Delta y_t^{10}>0,\\
\text{falling}, & \Delta y_t^{10}\leq0.
\end{cases}
\]

The default lookback is 13 weeks. The app also offers 4-, 8-, and
26-week windows.

\subsection{Supporting indicators}

The 2s10s curve slope is:

\[
\mathrm{Slope}_t=100\left(y_t^{10}-y_t^{2}\right)
\]

in basis points. A positive value indicates a conventional upward-sloping
curve.

The 10-year breakeven rate is displayed as a market measure of average
inflation compensation over the next decade. Neither supporting indicator
changes the regime classification in version 0.1.

\section{Four-Regime Framework}

\begin{longtable}{p{0.17\textwidth}p{0.13\textwidth}p{0.25\textwidth}p{0.27\textwidth}}
\toprule
\textbf{Profits} & \textbf{10Y yield} & \textbf{Regime} &
\textbf{Interpretation} \\
\midrule
\endhead
Accelerating & Rising &
Structural growth / higher cost of capital &
Nominal growth is stronger, but discount rates and the cost of capital are
also rising. \\
Accelerating & Falling &
Risk-on easing &
Profit momentum is improving while financial conditions become more
supportive. \\
Decelerating & Rising &
Margin pressure / stagflation risk &
Profit momentum is weakening while discount rates remain under upward
pressure. \\
Decelerating & Falling &
Cyclical slowdown &
Both profit momentum and long-term yields point toward slower nominal growth. \\
\bottomrule
\end{longtable}

The classification rule is intentionally deterministic:

\begin{lstlisting}[style=code]
def classify_regime(profits_accelerating, yields_rising):
    if profits_accelerating and yields_rising:
        return "Structural growth / higher cost of capital"
    if profits_accelerating:
        return "Risk-on easing"
    if yields_rising:
        return "Margin pressure / stagflation risk"
    return "Cyclical slowdown"
\end{lstlisting}

\section{Processing Pipeline}

The calculation runs in the following order:

\begin{enumerate}
  \item Download the four FRED CSV files.
  \item Parse dates and convert values to numeric observations.
  \item Calculate quarterly year-over-year profit growth.
  \item Calculate the quarterly change in profit growth.
  \item Resample the daily rate series to weekly Friday observations.
  \item Carry the latest quarterly profit signal into the weekly table.
  \item Calculate the lookback change in the 10-year yield.
  \item Calculate the 2s10s slope.
  \item Classify every week into a regime.
  \item Display the current regime and historical series.
\end{enumerate}

\section{Application Design}

\subsection{Interface}

The Streamlit application contains:

\begin{itemize}
  \item a sidebar for chart history and yield lookback;
  \item a current-regime card with a short interpretation;
  \item five metric cards;
  \item a dual-axis historical chart;
  \item colored background bands showing past regimes;
  \item a table of recent weekly observations;
  \item a CSV download; and
  \item an expandable methodology note.
\end{itemize}

\subsection{Chart}

The chart places year-over-year corporate profit growth on the left axis and
the 10-year Treasury yield on the right axis. Background color indicates the
classified regime.

The two axes have different units and should not be compared by vertical
distance. The chart is intended to show direction, turning points, and regime
transitions.

\subsection{Caching}

Downloaded data are cached for one hour:

\begin{lstlisting}[style=code]
@st.cache_data(ttl=3600, show_spinner=False)
def get_history(lookback_weeks):
    data = load_market_data()
    return build_history(...)
\end{lstlisting}

Caching reduces repeated network requests when a user changes a display
control.

\section{Repository Structure}

\begin{verbatim}
profit-yield-regime-monitor/
|-- app.py
|-- regime_monitor/
|   |-- __init__.py
|   |-- data.py
|   |-- event_study.py
|   +-- signals.py
|-- tests/
|   +-- test_signals.py
|-- .streamlit/config.toml
|-- .github/workflows/tests.yml
|-- requirements.txt
|-- pyproject.toml
|-- README.md
+-- documentation.tex
\end{verbatim}

\texttt{data.py} handles downloads, while \texttt{signals.py} handles the
classification. The event calculations live in
\texttt{event\_study.py}. Presentation stays in \texttt{app.py}. This
separation allows the analytical logic to be tested without running Streamlit
or making a network request.

\section{Running the Project}

\begin{lstlisting}[style=code,language=bash,numbers=none]
git clone <repository-url>
cd profit-yield-regime-monitor
python -m venv .venv

# macOS or Linux
source .venv/bin/activate

# Windows PowerShell
.venv\Scripts\Activate.ps1

pip install -r requirements.txt
streamlit run app.py
\end{lstlisting}

Streamlit normally opens at \url{http://localhost:8501}.

\section{Testing}

Run:

\begin{lstlisting}[style=code,language=bash,numbers=none]
python -m unittest discover -s tests -v
\end{lstlisting}

The tests cover:

\begin{description}[style=nextline]
\item[Four regime combinations]
Confirms that each pair of Boolean signals maps to the expected regime.

\item[Signal calculations]
Uses synthetic quarterly profits and daily rates to check positive yield
change, curve slope, and the resulting regime.

\item[Empty input]
Confirms that requesting a snapshot from an empty history raises an error.

\item[Event dates and summaries]
Checks that only regime-entry dates become events and that the grouped sample
counts agree with the event-level data.
\end{description}

GitHub Actions runs the same tests on every push and pull request.

\section{Interpreting the Output}

\subsection{Structural growth / higher cost of capital}

This regime does not mean that higher yields are harmless for equities. It
means the profit and yield signals may share a common source, such as stronger
nominal activity, supply constraints, or investment demand. Equity outcomes
still depend on whether earnings growth is sufficient to offset higher
discount rates.

\subsection{Risk-on easing}

Improving profit momentum and falling yields are jointly supportive. The
classification does not distinguish falling yields caused by lower inflation
from those caused by concern about future growth. Breakevens and the yield
curve provide additional context.

\subsection{Margin pressure / stagflation risk}

This is the least favorable combination in this classification. Slower profit growth
reduces fundamental support while rising yields increase discount rates. The
label describes a signal combination, not a formal stagflation forecast.

\subsection{Cyclical slowdown}

Falling yields can support valuations, but decelerating profits point to weaker
activity. The net asset-price effect is not predetermined.

\section{Event Study}

The dashboard includes a short event study to check what followed each regime
in the available ETF history. An event occurs only when the regime label
changes. This avoids treating every week in a long regime as an independent
observation.

For asset $j$, event date $t$, and a 13-week horizon, the forward price return is

\[
R_{j,t}^{(13)} = 100\left(\frac{P_{j,t+13}}{P_{j,t}}-1\right).
\]

The two assets are SPY, used as a U.S. equity proxy, and TLT, used as a
long-duration Treasury proxy. Daily adjusted closes are downloaded from Yahoo
Finance and resampled to Friday observations. For each regime and asset, the app
reports the number of events, median return, mean return, standard deviation,
and share of positive returns.

The chart uses the median because a few large market moves can pull the mean
away from a typical event. The event count is shown with the result; small
groups should not be read as stable estimates.

This exercise is descriptive. It does not include transaction costs and relies
on the latest vintage of the revised profit series. Regimes may also overlap in their 13-week outcome
windows when signals change quickly. The results are a historical check on the
classification, not evidence of a trading strategy.

\section{Limitations}

\begin{enumerate}
  \item \textbf{The profit proxy is backward-looking.} BEA profits do not
  measure analyst revisions.
  \item \textbf{Mixed frequencies create stale signals.} The profit signal
  changes quarterly while yields change daily.
  \item \textbf{Data are revised.} Historical classifications may change after
  BEA revisions.
  \item \textbf{The thresholds are binary.} A one-basis-point yield change and
  a large yield move receive the same direction label.
  \item \textbf{The lookback is selected by the user.} Different windows can
  produce different regimes.
  \item \textbf{No asset-return forecast is estimated.} The labels are market
  descriptions, not trading signals.
  \item \textbf{The model is U.S.-focused.} Global profit, currency, and
  policy differences are not included.
  \item \textbf{The event study uses ETF adjusted closes.} The sample begins
  with available ETF data and depends on a third-party price source.
  \item \textbf{Event windows can overlap.} Closely spaced regime changes do
  not produce statistically independent observations.
\end{enumerate}

\section{Possible Extensions}

The project can remain compact while adding one extension at a time:

\begin{itemize}
  \item add neutral zones around zero to reduce frequent switching;
  \item report regime persistence and transition counts;
  \item compare 4-, 13-, and 26-week event horizons;
  \item add a real-yield series;
  \item compare BEA profits with a user-supplied analyst-estimate CSV; or
  \item record data vintages to study the effect of revisions.
\end{itemize}

\section{Data Sources}

\begin{itemize}
  \item U.S. Bureau of Economic Analysis, Corporate Profits:
  \href{https://www.bea.gov/data/income-saving/corporate-profits}{BEA corporate profits page}
  \item FRED Corporate Profits After Tax:
  \href{https://fred.stlouisfed.org/series/CP}{FRED series CP}
  \item FRED 10-Year Treasury Yield:
  \href{https://fred.stlouisfed.org/series/DGS10}{FRED series DGS10}
  \item FRED 2-Year Treasury Yield:
  \href{https://fred.stlouisfed.org/series/DGS2}{FRED series DGS2}
  \item FRED 10-Year Breakeven Inflation:
  \href{https://fred.stlouisfed.org/series/T10YIE}{FRED series T10YIE}
  \item Yahoo Finance Historical Prices:
  \href{https://finance.yahoo.com/}{Yahoo Finance}
\end{itemize}

\section{Conclusion}

The monitor is deliberately narrow. It takes one question from market
commentary---how profits and yields are moving relative to one another---and
turns it into two explicit signals. The result is easy to reproduce and easy
to disagree with. Every judgment is visible, from the earnings proxy to the
13-week yield window.

\vfill
\begin{center}
\fcolorbox{acid}{paper}{
\parbox{0.87\textwidth}{
\centering
\textbf{Disclaimer}\\
For research only. Data may be revised. This project does
not provide investment advice.
}}
\end{center}

\end{document}
